\section{Experimental Evaluation and Discussion}
\label{sec:results}

\subsection{Protocol and Metrics}

The student is trained on \(0\)--\(80\%\) of each recording and validated on
\(80\)--\(90\%\). The frozen SHARP classifier was trained on S1a/S1b/S1c over
\(0\)--\(60\%\) and selected on \(60\)--\(80\%\); no classifier weights are
updated. Both are evaluated on the same disjoint \(90\)--\(100\%\) interval,
with stride 30 and 718 windows from the five classifier activities. This is
generalization to unseen temporal portions of known recordings, not to unseen
recordings, people, or environments. Each method receives identical raw
windows. The recording-mean centering expected by the classifier is estimated
independently for every representation and does not use labels.

We report full-map and off-center MSE against windowed targets from the official
precomputed SHARP Doppler traces. A frame is active when its maximum outside
bins 45--55 exceeds 0.2;
precision, recall, and F1 compare predicted and target support. The motion-mass
ratio is total predicted divided by target off-center excess power, with one
ideal. Frozen classifier accuracy is the primary task metric.

\subsection{RQ1: Can the Student Learn the Mapping?}

\textbf{Memorization sanity check.}
Before full-data training, two students were optimized repeatedly on the same
16 motion-rich PI windows. The 30-subcarrier model reached
\(3.881\times10^{-5}\) MSE, while the 242-subcarrier model reached
\(4.194\times10^{-5}\)(\figref{fig:diagnostics}(a)). Both reproduced the selected peaks. Thus, the end-to-end training path is functioning
correctly, and the model has sufficient capacity to memorize the selected local
examples, although this test says nothing about generalization.

\textbf{Full-data behavior.}
With official precomputed SHARP Doppler supervision, the final model behaves
very differently from the preliminary students. Motion-aware training loss decreases from
0.03397 to 0.01264 over all 50 epochs. All-scenario validation decreases from
0.02800 to 0.01415, while S1a/S1b/S1c validation decreases from 0.02506 to
0.01326 and reaches its minimum, 0.01212, at epoch 45 (\figref{fig:diagnostics}(b)). Both validation curves
continue improving, so there is no observed deterioration on the held-out
temporal intervals. Because training uses motion-balanced sampling while
validation retains the natural distribution, their absolute loss gap is not
interpreted.

\begin{figure*}[t]
    \centering
    \includegraphics[width=\textwidth]{diagnostics.pdf}
    \caption{Optimization diagnostics. (a) Both 30- and 242-subcarrier models
    memorize 16 fixed windows. (b) With official precomputed SHARP Doppler
    supervision, full-data
    train and validation objectives improve throughout training; the dashed
    line marks the selected epoch-45 checkpoint.}
    \label{fig:diagnostics}
\end{figure*}

The held-out example in \figref{fig:qualitative} shows that the model has learned an input-dependent mapping rather than producing a fixed average Doppler pattern. The student places off-center responses at the target event times. Its main error is now spatial smoothing:
thin target peaks become wider lobes, and nearby temporal events merge. The
complete corrected configuration avoids the center-only shortcut in this
example, but does not preserve all fine motion geometry.

\begin{figure*}[t]
    \centering
    \includegraphics[width=\textwidth]{qualitative.pdf}
    \caption{Held-out S1a\_C window 16920--17260 with identical color limits.
    The full-resolution student localizes the principal target events but
    broadens them; affine sanitization followed by the fixed STFT retains
    sharper time--Doppler structure.}
    \label{fig:qualitative}
\end{figure*}

\subsection{RQ2: Do Reconstructed Doppler Traces Preserve Downstream HAR
Performance Under a Frozen Classifier?}

\tabref{tab:main-results} gives the matched comparison. The antenna-shared
frequency--time U-Net reaches \(63.09\%\), a 19.78-point improvement over the
antenna-shared temporal U-Net with a coarse spectral head and well above raw
STFT. The complete corrected architecture and training
configuration therefore improves both map and downstream metrics.

\begin{table*}[t]
\centering
\caption{Held-out S1a/S1b/S1c comparison over the same 718 windows. All map
metrics use windowed targets from the official precomputed SHARP Doppler
traces.}
\label{tab:main-results}
\footnotesize
\setlength{\tabcolsep}{5pt}
\begin{tabular}{@{}lrrrrr@{}}
\hline
Representation & Accuracy (\%) & Balanced (\%) & Off-center MSE & Active F1 &
Mass ratio \\
\hline
Official precomputed SHARP Doppler & \textbf{91.36} & 91.74 & 0 & 1.000 & 1.00 \\
Raw fixed STFT & 21.73 & 20.00 & 0.128148 & 0.357 & 33.28 \\
Direct affine sanitization + fixed STFT & \textbf{91.36} & \textbf{91.80} & \textbf{0.000155} & \textbf{0.943} & \textbf{0.98} \\
Full-resolution frequency--time student & 63.09 & 58.86 & 0.001062 & 0.758 & 1.21 \\
Shared temporal student + coarse spectral head & 43.31 & 40.05 & 0.002758 & 0.505 & 2.22 \\
\hline
\end{tabular}
\end{table*}

The \(63.09\%\) aggregate accuracy does not imply classifier-equivalent or
uniform recognition: it is helped by \(100\%\) empty-room, \(99.36\%\) still,
and \(91.67\%\) running accuracy. Walking and jumping reach only \(1.30\%\) and
\(1.98\%\), respectively; 152 of 154 walking windows and 92 of 101 jumping
windows are classified as running. Balanced accuracy is consequently only
\(58.86\%\). The student has learned whether and approximately where motion
occurs, as supported by F1 \(0.758\) and mass ratio 1.21, but smoothing removes
the spectral and temporal details needed to distinguish dynamic activities.

In this provenance-matched experiment, off-center MSE and support F1 broadly
rank the methods in the same order as classifier accuracy. They remain
insufficient by themselves: active support treats a broad running-like lobe as
successful even when it erases class-discriminative shape. Reconstruction and
task metrics must therefore be reported together.

\subsection{RQ3: What Does a Simplified Phase-Sanitization Baseline Reveal?}

We include a non-learned baseline to test whether explicitly encoding a known
phase-error structure is sufficient, rather than asking the student to discover
it. The procedure is adapted from the affine phase-sanitization stage in
SHARP's public pipeline \cite{meneghello2023sharp}, consistent with established
CSI phase-error models \cite{ratnam2024optimal}. It applies only this correction
directly to raw CSI, omits SHARP's sparse path estimation, and then uses the
fixed Doppler transform.

This baseline matches the classifier accuracy of the official precomputed
SHARP Doppler at \(91.36\%\). Their individual predictions agree on \(98.33\%\) of
windows. Thus, the methods make different mistakes on a small number of samples, and these differences cancel out, resulting in the same overall accuracy. The affine
method obtains off-center MSE \(1.55\times10^{-4}\), F1 \(0.943\), and
motion-mass ratio 0.982.

This establishes that the local raw windows contain enough information for a
high-utility representation. It also highlights the main remaining neural error:
the direct model learns the approximate timing and location of motion despite unwanted phase noise. However, its generic decoder often replaces SHARP’s sharp, sample-specific peaks with smoother average patterns.
Under this setup, explicit affine phase modeling is highly effective and leaves
a substantially smaller gap than the corrected generic neural decoder.

\subsection{Limitations and Answer to the RQ}

The experiments use one seed, one frozen classifier, and 92 aligned AR
recordings paired with official precomputed SHARP Doppler traces; five candidate pairs are
excluded for temporal mismatch. The classifier test measures compatibility
with SHARP's established interface, not
the performance of a newly trained classifier adapted to student maps. 

The answer to the primary question is therefore mixed but operationally
negative. Direct regression can learn a meaningful approximation: the corrected
model generalizes event support and reaches \(63.09\%\), so the task is not
infeasible. It does not preserve enough fine structure to replace SHARP at
\(91.36\%\), especially for walking and jumping. The most promising next step is a phase-aware hybrid model that first cleans the CSI and then learns only the remaining correction or transformation.
